% ID: SEL-D0007-BD-E07
% ESTADO: cotejada. Segunda pasada visual completa y borrador compilado revisados.
% Fuente: Selectividad/matematicas_aplicadas/ORD_TITULAR_MATEMÁTICAS APLICADAS CCSS-JUNIO-IMPRESO.pdf
% SHA-256: c42a605501a277fe7194b964ad92ad66047be9a62dc254fa2219d756d3d71317
% Página: 2; bloque D; ejercicio original 7. Un único 1.5 para a.1 y a.2 juntos.

\begin{enumerate}[label=\alph*)]
  \item \textbf{(1.5 puntos)} En una Escuela Politécnica hay matriculados
  en el último curso 60 estudiantes de Ingeniería Eléctrica, 40 de Ingeniería
  Informática, 30 de Ingeniería Civil, 50 de Ingeniería Mecánica y 20 de
  Ingeniería Aeronáutica. Se quiere hacer una encuesta al 20\,\% de estos
  estudiantes, de manera proporcional al número de matriculados en cada
  titulación.
  \begin{enumerate}[label=\arabic*.]
    \item \textquestiondown{}Qué tipo de muestreo se debe emplear?
    \item \textquestiondown{}Cuántos alumnos debe haber en la muestra y
    cuántos de cada titulación?
  \end{enumerate}
  \item \textbf{(1 punto)} Dada la población $\{a,10,12,11,18\}$,
  \textquestiondown{}cuánto debe valer $a$, sabiendo que la media de las
  medias muestrales de tamaño 3, obtenidas mediante muestreo aleatorio
  simple, es 13.2?
\end{enumerate}
